\uniterablesection{Вступ}

Звіт відображає результати передкваліфікаційної практики, метою якої було формалізувати, дослідити та верифікувати вебзастосунок для тривимірного моделювання внутрішньобудинкової електромережі приміщень з автоматичним розрахунком довжини проводки та витрат оздоблювальних матеріалів. Застосунок дозволяє побудувати квартиру з кількох кімнат, розмістити в ній двері, арки, розетки, вимикачі, щиток, розподільні коробки та світильники, прокласти між ними проводи з ортогональною трасою, а потім отримати кількісну специфікацію: площі стін, кількість листів гіпсокартону (ГКЛ), профілів і плінтусів, довжину кожного проводу. Окремий модуль будує розкладку листів ГКЛ по стінах і стелі з повторним використанням відрізків.

\textbf{Актуальність} теми визначається такими чинниками. По-перше, планування внутрішньої електропроводки та закупівля матеріалів на практиці зводяться до окремих ручних підрахунків, які виконують за планом приміщення; нормативні вимоги до проєктування електроустановок житлових будинків зосереджено в ДБН~В.2.5-23~\citecustom{6}, але перелік необхідних довжин кабелю та кількості матеріалів з них автоматично не випливає. По-друге, універсальні редактори планування приміщень (Sweet Home 3D, Planner 5D, RoomSketcher) орієнтовані на дизайн інтер'єру: електричні символи в RoomSketcher доступні лише в платних тарифах~\citecustom{8}, а проводка як просторовий маршрут з розрахованою довжиною в описах їхніх можливостей не заявлена~\citecustom{7}--\citecustom{9}. По-третє, онлайн-калькулятори гіпсокартону визначають кількість листів за площею з відсотком запасу й не будують розкладку листів~\citecustom{10}, \citecustom{11}, хоча задача розкрою листового матеріалу є добре дослідженою оптимізаційною задачею~\citecustom{1}--\citecustom{3}. Виявлена прогалина полягає у відсутності легкого вебзасобу, який поєднує в одній моделі геометрію приміщення, електрообладнання, траси проводів і розрахунок матеріалів із повторним використанням відрізків.

\textbf{Мета практики} --- проаналізувати предметну область, сформулювати специфікацію вимог, формалізувати архітектуру, моделі та алгоритми вебзастосунку, а також експериментально перевірити коректність його розрахунків.

Для досягнення мети поставлено такі \textbf{завдання}:
\begin{enumerate}[noitemsep, topsep=0.3em, leftmargin=*, label=\arabic*.]
    \item проаналізувати проблематику планування внутрішньої електромережі приміщень, розрахунку оздоблювальних матеріалів і розкрою листових матеріалів;
    \item виконати огляд програмних аналогів і визначити нішу застосунку;
    \item розробити специфікацію вимог за ISO/IEC/IEEE 29148:2018 та модель варіантів використання;
    \item обґрунтувати архітектурний стиль, описати компоненти, модель даних і геометричну модель приміщення;
    \item формалізувати алгоритми трасування проводів, розрахунків матеріалів і розкрою ГКЛ;
    \item верифікувати розрахунки незалежним перерахунком, виміряти продуктивність розкрою та окреслити обмеження поточної версії.
\end{enumerate}

\textbf{Об'єкт дослідження} --- процес планування розміщення електрообладнання та проводки в приміщеннях і підготовки специфікації матеріалів. \textbf{Предмет дослідження} --- методи й програмні засоби тривимірного моделювання приміщень, ортогонального трасування проводів і розкрою листових матеріалів у форматі вебзастосунку.

\textbf{Методи дослідження}: системний аналіз предметної області, інженерія вимог (ISO/IEC/IEEE 29148), об'єктно-орієнтоване моделювання (UML), архітектурне моделювання за нотацією C4, порівняльний аналіз аналогів, а також експериментальна перевірка: незалежний перерахунок контрольної сцени окремим скриптом і вимірювання часу виконання алгоритму розкрою.

\textbf{Практичне значення} результатів полягає в тому, що описаний вебзастосунок працює в браузері без сервера та бази даних, зберігає проєкт у файлі JSON і на контрольній сцені з двох кімнат дає специфікацію, яка збігається з незалежним ручним розрахунком. Режим розкрою «Економія» зменшив кількість листів ГКЛ із 33 (оцінка за площею) до 30, що відповідає нижній межі за площею поверхні. Нормативно-обчислювальна частина електротехнічних розрахунків (навантаження, переріз, падіння напруги) у поточній версії не реалізована й визначена як наступний етап кваліфікаційної роботи.

\textbf{Структура звіту.} Звіт складається з календарного плану, вступу, двох розділів, висновків, списку використаних джерел. У розділі~I проаналізовано предметну область, виконано огляд аналогів і сформульовано вимоги. У розділі~II описано архітектуру, модель даних, геометричну модель, алгоритми трасування, розрахунків і розкрою, інтерфейс користувача та результати перевірки.
